% =============================================================================
% ABSTRACT (01_abstract.tex)
% =============================================================================

\section{Abstract}
\addcontentsline{toc}{section}{Abstract}

\begin{comment}
Im Rahmen meiner Bachelorarbeit wird das Prüfverfahren für MCC-Felder der Schaltgerätekombination TopDrawPlus (Rolf Janssen GmbH) durch eine automatisierte Stromregelung optimiert. Ziel ist es, den Bemessungsbetriebsstrom stufenlos konstant zu halten, um thermische Drifts im Prüfstand ohne Verletzung der normativen Vorgaben nach DIN EN 61439 zu kompensieren. Bisherige manuelle Nachregulierungen führten zu hohen Zeitaufwänden und einer eingeschränkten Reproduzierbarkeit der Messergebnisse

Im ersten Schritt werden die physikalischen Ursachen der Instabilitäten systematisch untersucht. Hierbei stehen die temperaturabhängige Widerstandszunahme der Strombahnen, die thermische Kopplung parallel betriebener Prüflinge sowie elektromagnetische Wechselwirkungen bei hohen Strömen im Fokus. Auf Basis dieser experimentellen Daten werden mathematische Modelle zur Beschreibung der thermo-elektrischen Systemdynamik entwickelt und validiert.

Darauf aufbauend werden dezentrale Regelungskonzepte entworfen, die eine individuelle Kompensation gegenseitiger Beeinflussungen ermöglichen. Diese Ansätze integrieren moderne Regelalgorithmen mit einer robusten Hardwarearchitektur, bestehend aus präziser Sensorik, Leistungselektronik und einer übergeordneten Steuerungsebene.

Mittels einer Nutzwertanalyse werden verschiedene Lösungsvarianten hinsichtlich Normkonformität, thermischer Dauerlaststabilität sowie technischer und wirtschaftlicher Machbarkeit bewertet. Für die ermittelte Vorzugsvariante erfolgt die detaillierte Ausarbeitung der Hard- und Softwarearchitektur. Dies beinhaltet die Implementierung einer Steuerungslogik zur automatisierten Überwachung normativer Abbruchkriterien sowie zur Sicherstellung der geforderten Temperaturstabilität über die gesamte Prüfdauer.

Die entwickelte Automatisierungslösung ersetzt das manuelle Nachregeln, steigert die Effizienz der Prüfzyklen und minimiert prozessbedingte Fehlerquellen deutlich. Abschließend wird die Lösung in die Prüffeld-Dokumentation überführt und dient als technische Grundlage für zukünftige Bauartnachweise.

\end{comment}

\textcolor{red}{\textbf{Achtung! Überarbeiten}}